\appendix
\section{Export and Runtime Details}
\label{app:export}

The original generation scripts load RGB TIFFs and call the pretrained
image-conditioned pipeline with seed 1. Captions are saved separately, not passed
into generation. The first mesh and Gaussian output are passed to the standard
GLB exporter with \texttt{simplify=0.95} and \texttt{texture\_size=1024}; the
Gaussian is also exported as PLY. The new ablation requires only GLB export and
holds those export arguments fixed. The simplification argument is reported
literally rather than asserting a retained-face fraction across software versions.

The historical full-batch script skips existing GLBs. It records 433 generation
entries for 438 archived meshes; skipped outputs do not acquire fresh records.
CUDA synchronization brackets the pipeline call, so the timer excludes initial
model loading, TIFF loading, and later export. Peak PyTorch-allocated memory is
divided by $1024^3$ (GiB) and includes model allocations. It is not total GPU usage
or an end-to-end deployment cost. The ablation logs setup, pipeline, and export
separately; attention backends must be consistent across compared runs.

\section{Preprocessing and Diagnostic Definitions}
The adaptation preserves aspect ratio on a white canvas, thresholds ink with
Otsu's rule, and applies an unsigned Euclidean distance transform to non-ink
pixels. There is no stroke opening or speck deletion in this stage. The Gaussian
map is quantized to RGB and used only as the generator input. The evaluation
reference remains the original reviewed crop for all branches.

The separate legacy evaluator thresholds grayscale, chooses polarity from border
pixels, applies opening and closing with a $3\times3$ kernel, and removes
components smaller than $\max(16,\lfloor0.0002HW\rfloor)$. A foreground box is
expanded by eight percent and centered after aspect-preserving resizing. This
procedure can remove sparse ink. Re-thresholding saved aligned references is a
diagnostic rather than a reconstruction of the pre-resize mask.

Recomputed mesh connectivity uses native indices. The closed-edge test requires
two incident triangles per undirected edge and does not test self-intersection
or every manifold condition. Degenerate triangles use a cross-product norm
threshold of $10^{-12}$ in native coordinates. These definitions differ from
vertex-welded or face-connected mesh analyses and should accompany comparisons.

